\section*{Capítulo \ref{ch:g7:volcanoes} --- Los volcanes}

\begin{solution}{exo:g7:volcanoes:1}
(a) \omterm{def:g7:volcanoes:vent}{cráter}; (b) \omterm{def:g7:volcanoes:magma}{magma}; (c) \omterm{def:g7:volcanoes:vent}{boca eruptiva}; (d) \omterm{def:g4:volcanoes-earthquakes-first-look:eruption}{lava}; (e) \omterm{def:g7:volcanoes:magma-chamber}{cámara magmática}.
\end{solution}

\begin{solution}{exo:g7:volcanoes:2}
(a) Efusiva. (b) Explosiva. (c) Efusiva: una fuente de \omterm{def:g4:volcanoes-earthquakes-first-look:eruption}{lava} lanza gotas
incandescentes que vuelven a caer y corren; el \omterm{def:g7:volcanoes:magma}{magma} no se convierte en
ceniza. (d) Explosiva: es un \omterm{def:g7:volcanoes:pyroclastic-flow}{flujo piroclástico}.
\end{solution}

\begin{solution}{exo:g7:volcanoes:3}
(a) Verdadero. (b) Falso: el gas volcánico más abundante es el vapor de
agua; el dióxido de carbono viene después. (c) Falso: un \omterm{def:g7:volcanoes:pyroclastic-flow}{flujo piroclástico}
es una avalancha de gas caliente, ceniza y trozos de \omterm{def:g3:stones-pebbles-sand:rock}{roca}. (d) Falso: están
en el interior de la \omterm{def:g7:plate-tectonics:plate}{placa} del Pacífico, a más de 3200~km del límite más
cercano, sobre un \omterm{def:g7:volcanoes:hotspot}{punto caliente}.
\end{solution}

\begin{solution}{exo:g7:volcanoes:4}
(a) De unos 1000 a 1200~grados. (b) De unos 600 a 800~grados. (c) Cerca de
la \omterm{def:g7:volcanoes:vent}{boca}, donde la \omterm{def:g4:volcanoes-earthquakes-first-look:eruption}{lava} es amarilla (de 1000 a 1200~grados) y no naranja (de
800 a 1000~grados): la \omterm{def:g4:volcanoes-earthquakes-first-look:eruption}{lava} se enfría a medida que se aleja de la \omterm{def:g7:volcanoes:vent}{boca}.
\end{solution}

\begin{solution}{exo:g7:volcanoes:5}
(a) $838 + 221 = 1059$ \omterm{def:g4:volcanoes-earthquakes-first-look:volcano}{volcanes}. (b) $1059 \div 1214 \approx 0.872$:
alrededor del \textbf{87~\%}. (c) $69 + 18 + 12 + 0.5 = 99.5$, alrededor
del 100~\%: la pequeña diferencia se debe al redondeo de cada porcentaje.
\end{solution}

\begin{solution}{exo:g7:volcanoes:6}
(a) $564 \div 660 \approx 0.855$: alrededor del \textbf{85~\%}. (b)
$660 - 564 - 62 - 33 = 1$: un cono empinado. (c) Los conos empinados los
construyen la \omterm{def:g4:volcanoes-earthquakes-first-look:eruption}{lava} viscosa y la ceniza, la marca de las \omterm{def:g7:volcanoes:eruption-types}{erupciones
explosivas}: a lo largo del Anillo de Fuego, donde una \omterm{def:g7:plate-tectonics:plate}{placa} se hunde bajo
otra, cabe esperar muchas \omterm{def:g7:volcanoes:eruption-types}{erupciones explosivas}, como las del monte Santa
Helena y del Pinatubo.
\end{solution}

\begin{solution}{exo:g7:volcanoes:7}
(a) $1170 - 1140 = 30$~grados. (b) $30 \div 10 = 3$~grados por kilómetro.
(c) En el túnel, la \omterm{def:g4:volcanoes-earthquakes-first-look:eruption}{lava} no toca el aire: el techo endurecido que tiene
encima conserva el calor, como la tapa de una olla. Al aire libre, su
superficie pierde calor enseguida.
\end{solution}

\begin{solution}{exo:g7:volcanoes:8}
(a) Kauai, por al menos $5.5 - 0.7 = 4.8$~millones de años. (b) Kauai: la
\omterm{def:g4:clouds-rain-snow:precipitation}{lluvia} y las olas llevan mucho más tiempo desgastándola. (c) La Isla
Grande, donde todavía sale \omterm{def:g4:volcanoes-earthquakes-first-look:eruption}{lava} nueva. (d) Hacia el noroeste: Kauai, que se
alejó primero del \omterm{def:g7:volcanoes:hotspot}{punto caliente}, está ahora al noroeste de él.
\end{solution}

\begin{solution}{exo:g7:volcanoes:9}
(a) Casi todos los \omterm{def:g4:volcanoes-earthquakes-first-look:volcano}{volcanes} están en \omterm{def:g7:plate-tectonics:plate-boundary}{límites de placas} (el 87~\% de la
lista): unos \omterm{def:g4:volcanoes-earthquakes-first-look:volcano}{volcanes} tan lejos de cualquier límite son una excepción.
(b) En las profundidades del \omterm{def:g6:inside-the-earth:layers}{manto} bajo Hawái, un \omterm{def:g7:volcanoes:hotspot}{punto caliente} se
mantiene caliente durante millones de años. Funde en parte la \omterm{def:g3:stones-pebbles-sand:rock}{roca} de
encima. El \omterm{def:g7:volcanoes:magma}{magma}, más ligero que la \omterm{def:g3:stones-pebbles-sand:rock}{roca} sólida que lo rodea, sube a través
de la \omterm{def:g7:plate-tectonics:plate}{placa} y entra en \omterm{def:g4:volcanoes-earthquakes-first-look:eruption}{erupción}. A medida que la \omterm{def:g7:plate-tectonics:plate}{placa} se mueve, crecen
\omterm{def:g4:volcanoes-earthquakes-first-look:volcano}{volcanes} nuevos sobre el \omterm{def:g7:volcanoes:hotspot}{punto caliente} y los viejos se van alejando. (c)
Puntos naranjas (``en el interior de una \omterm{def:g7:plate-tectonics:plate}{placa}'').
\end{solution}

\begin{solution}{exo:g7:volcanoes:10}
El primer \omterm{def:g7:volcanoes:magma}{magma} había perdido la mayor parte de su gas al subir, como una
botella de agua con gas que se ha dejado abierta o quieta: cuando llegó a
la superficie, ya no quedaba nada que lo hiciera estallar, y salió
lentamente. El \omterm{def:g7:volcanoes:magma}{magma} del 12 y del 15~de junio todavía conservaba su gas,
como la botella agitada: al bajar la presión cerca de la superficie, el gas
formó burbujas de golpe, que se hincharon y desgarraron el \omterm{def:g7:volcanoes:magma}{magma} viscoso en
ceniza. El primer \omterm{def:g7:volcanoes:magma}{magma} es el que había perdido su gas.
\end{solution}

\begin{solution}{exo:g7:volcanoes:11}
(a) $10 \div 80 = 0.125$~horas, y $0.125 \times 60 = 7.5$~minutos. (b)
$10 \div 130 \approx 0.077$~horas, unos $0.077 \times 60 \approx
4.6$~minutos, es decir, unos 5~minutos. (c) La familia necesita
20~minutos, pero el flujo llega en 5 a 8~minutos: cuando se ve el flujo, ya
es demasiado tarde. La familia solo puede estar a salvo si se marcha cuando
los científicos ven las señales de aviso, antes de la \omterm{def:g4:volcanoes-earthquakes-first-look:eruption}{erupción}.
\end{solution}

\begin{solution}{exo:g7:volcanoes:12}
(a) $111 \times 360 = 39\,960$: unos \textbf{40\,000~km}. (b)
$40\,000 \div 15 \approx 2667$: unos 2\,670~km al día. (c)
$2667 \div 24 \approx 111$: unos \textbf{110~km por hora}. (d) El viento:
la ceniza fina no tiene velocidad propia una vez que la \omterm{def:g4:volcanoes-earthquakes-first-look:eruption}{erupción} la ha
lanzado hacia arriba; flota muy alto en el aire y el viento la lleva,
mientras cae despacio.
\end{solution}

\begin{solution}{pb:g7:volcanoes:1}
\textbf{1.} Del 16 al 31~de marzo: $31 - 16 = 15$~días; después, los
30~días de abril y 18~días de mayo: $15 + 30 + 18 = 63$~días.

\textbf{2.} $10\,000 \div 62 \approx 161$: más de 160 pequeños \omterm{def:g4:volcanoes-earthquakes-first-look:earthquake}{terremotos}
al día, de media.

\textbf{3.} La hinchazón mostraba que el \omterm{def:g7:volcanoes:magma}{magma} había subido muy arriba
dentro de la \omterm{def:g1:landscapes-around-us:landform}{montaña} y empujaba la \omterm{def:g3:stones-pebbles-sand:rock}{roca} hacia fuera. Además, volvía la
ladera norte empinada e inestable, lista para deslizarse.

\textbf{4.} Un triángulo azul está donde una \omterm{def:g7:plate-tectonics:plate}{placa} se hunde bajo otra.
Cabe esperar un \omterm{def:g7:volcanoes:magma}{magma} viscoso y rico en gas, \omterm{def:g7:volcanoes:eruption-types}{erupciones explosivas} y un
cono empinado (un estratovolcán) --- justo lo que muestra la primera
fotografía.

\textbf{5.} $240 \div 60 = 4$~kilómetros por minuto.

\textbf{6.} La ladera hinchada retenía el \omterm{def:g7:volcanoes:magma}{magma}, a presión, como el tapón de
una botella. Cuando se deslizó, la presión bajó de golpe; el gas disuelto
en el \omterm{def:g7:volcanoes:magma}{magma} se convirtió en burbujas de golpe y salió disparado por la
ladera abierta --- de lado, por donde había desaparecido la tapa.

\textbf{7.} $27 \div 480 \approx 0.056$~horas, y $0.056 \times 60 \approx
3.4$~minutos. Como la explosión avanzó al menos a esa velocidad, tardó
como mucho unos 3.4~minutos.

\textbf{8.} $8 \div 130 \approx 0.062$~horas, y $0.062 \times 60 \approx
3.7$: unos \textbf{4~minutos}.

\textbf{9.} $24 \div 8.8 \approx 2.7$: la columna subió unas 2.7 veces más
alto que el Everest.

\textbf{10.} Efusivas: la \omterm{def:g4:volcanoes-earthquakes-first-look:eruption}{lava} viscosa salió sin explotar y se amontonó
sobre la \omterm{def:g7:volcanoes:vent}{boca}. Como la botella que se había dejado quieta, este \omterm{def:g7:volcanoes:magma}{magma} ya
había perdido la mayor parte de su gas.

\textbf{11.} Antes, la \omterm{def:g1:landscapes-around-us:landform}{montaña} era un cono liso y puntiagudo cubierto de
\omterm{def:g4:clouds-rain-snow:precipitation}{nieve}, sobre bosques verdes. Después, su cima y su ladera norte han
desaparecido: en su lugar hay un enorme \omterm{def:g7:volcanoes:vent}{cráter} abierto hacia nosotros, y las
laderas y el \omterm{def:g1:landscapes-around-us:landform}{valle} de delante están grises de ceniza y de \omterm{def:g3:stones-pebbles-sand:rock}{roca} rota.

\textbf{12.} $2950 - 2549 = 401$: la \omterm{def:g1:landscapes-around-us:landform}{montaña} perdió unos
\textbf{401~m} de altura, unos 400~metros, como se contó en el curso~4.
\end{solution}
